\chapter{研究背景与问题界定}

\section{研究场景}

博士后研究报告兼具“科研总结”和“正式归档”双重属性。与一般课程论文相比，它不仅需要完整叙述研究目标、技术路线、结果与局限，还需要在封面、题名页、摘要、目录、插图和附表清单、参考文献、附录和结尾部分上满足相对严格的编排要求。本公开项目因此不追求某一具体课题的学术深度，而是优先验证结构完整性和工程稳定性。

\section{选题设定}

为了遵守本仓库毕业论文示例的统一约定，本示例选择“东亚流行文化语境下佐佐木希影像传播与审美接受研究”作为承载主题。该主题具备三个优点：一是全部内容可以基于公开文化现象和占位材料展开；二是适合演示图像、文本、社交媒体与跨区域接受等多类材料；三是便于覆盖图表、附录和结尾材料的完整链路。

\section{报告目标}

本示例聚焦三个目标：

\begin{enumerate}
  \item 给出一套可直接复用的中国科学院博士后出站报告项目骨架；
  \item 将《博士后研究报告编写规则》中的关键前置页要求，尤其是封面字段，转换为稳定的 LaTeX 结构；
  \item 在不改动 \texttt{bensz-thesis} 公共包实现的前提下，证明现有构建链路可以稳定承载该场景。
\end{enumerate}
